\subsection{Coprime modular covers with the old gentle conjugator retained}

\begin{statement}{D.44}{Proved}{Exact modular descent and gentle normalization over a coprime base}
Let
\(\widehat B=C_2\amalg C_3=\langle s,r\mid s^2=r^3=1\rangle\),
let
\(F=\ker(\widehat B\to\mathfrak S_3)\),
and use
\[
 \tau=sr,\qquad x=\tau^{-2},\qquad y=sxs,\qquad z=(xy)^{-1}.
 \tag{D.44.1}
\]
Let \(q:\widehat B\twoheadrightarrow\Gamma\) be an actual
pro-prime-to-\(p\) quotient retaining the displayed
\(\mathfrak S_3\) quotient, and put \(K=\ker q\).
Thus \(p\geq5\).
Suppose that an old automorphism
\(\alpha_0\in\operatorname{Aut}(\Gamma)\) satisfies
the following exact multiplier-one conditions:
it fixes \(\tau\), sends \(s,r\) to conjugates of
the same positive generators, and
\[
             f_0=s\alpha_0(s)
                   \ \in\ \operatorname{im}(F'\to\Gamma).
 \tag{D.44.2}
\]
Then on the actual relative cover
\[
                    G=\widehat B/O^p(K)
 \tag{D.44.3}
\]
there is an automorphism \(\sigma\) inducing exactly
\(\alpha_0\), fixing \(\tau\), and having
\[
 \begin{aligned}
 f=s\sigma(s)&\in\operatorname{im}(F'\to G),&
 q(f)&=f_0,\\
 \sigma(x)&=x,&
 \sigma(y)&=f^{-1}yf.
 \end{aligned}
 \tag{D.44.4}
\]
It satisfies the exact gentle equations
\[
       f(y,x)f(x,y)=1,\qquad
       f(z,x)f(y,z)f(x,y)=1
 \tag{D.44.5}
\]
in the actual pure quotient \(H=\operatorname{im}(F\to G)\).
Here evaluations use any one preimage of \(f\) in \(F'\);
their values are independent of that choice at the displayed
arguments.
Its pure outer automorphism centralizes the actual
\(\mathfrak S_3\)-outer action, by the explicit
comparison below.
The finite \(\sigma\)-stable quotient maps from this one \(G\) give a
compatible family of generating gentle shadows retaining
the literal old \(f_0\), and hence its old outer coordinates.

The torsion-conjugacy hypothesis in the statement holds
automatically when the pure kernel
\(\Gamma_F=\ker(\Gamma\to\mathfrak S_3)\)
is pro-\(\ell\) for \(\ell\geq5\), and the old
multiplier-one gentle equations define an automorphism
of \(\Gamma\).
In particular it holds for all the D.7--D.8
weight-nine gentle automorphisms at the actual
class-ten exponent-\(101\) quotient.
Each of them therefore has these complete coprime
relative lifts for every \(p\notin\{2,3,101\}\).
Its original three-strand shadow, including the
central coordinate and the given \(f_0\), is retained
by the actual common-source comparison below.
\end{statement}

\begin{proof}
We retain the indicated source generators throughout the construction.
\paragraph{The conventions and the actual pure free group.}
There is a surjection from the discrete free product
\(C_2*C_3\) to \(\mathfrak S_3\), with \(s\) a
transposition and \(r\) a three-cycle.
The transversal
\[
                      1,\ s,\ r,\ r^2,\ sr,\ sr^2
 \tag{D.44.6}
\]
gives twelve Schreier generators, one for each
transversal element and each of \(s,r\).
The only nonidentity reduced generators are
\[
 \begin{array}{c|c|c}
 \text{transversal element}&\text{letter}&
                          \text{Schreier generator}\\ \hline
 r&s&y^{-1}\\
 r^2&s&x\\
 sr&s&x^{-1}\\
 sr^2&s&y .
 \end{array}
 \tag{D.44.7}
\]
These are literal free-product equalities with
\(x,y\) from \textup{(D.44.1)}.
The kernel intersects every conjugate of the
two finite factors trivially; it acts freely
on the free-product tree.
Its quotient graph has six edges and five
vertices, hence free rank two.
The Schreier expressions in
\textup{(D.44.7)} show that \(x,y\) generate
that kernel and hence form a free basis.
Completion of this finite-index comparison
identifies the actual \(F\) with the free
profinite group on those same two words.

Direct free-product multiplication gives
\[
 \begin{array}{c|ccc}
       &x&y&z\\ \hline
 r(-)r^{-1}&y&z&x\\
 s(-)s&y&x&x^{-1}zx .
 \end{array}
 \tag{D.44.8}
\]
The cusp denoted \(t=r^{-1}s\) in the inverse
convention is \(\tau^{-1}\).
Preserving the conjugacy class of either
positive element is equivalent to preserving
the class of its inverse, by the actual
inversion equation.
For the pants source on \(a=s,b=r^{-1}\),
its third positive generator is
\[
             (ab)^{-1}=rs=s\tau s .
 \tag{D.44.9}
\]
Thus the pants source is not assigned the
word \(\tau\) as its third generator without
the conjugation displayed here.

\paragraph{The residual quotient and the chosen cusp conjugator.}
Let
\(\mathfrak F=\widehat F(a,b)\) be the free
profinite pants group and let
\(\pi:\mathfrak F\twoheadrightarrow\widehat B\)
send \(a\mapsto s,\ b\mapsto r^{-1}\).
Its kernel is the closed normal subgroup
\(\mathcal R\) generated by \(a^2,b^3\):
both quotients have exactly the same maps
to every finite group, by the indicated
generator relations.
Put \(K_{\mathfrak F}=\pi^{-1}(K)\).
The restriction
\(K_{\mathfrak F}\twoheadrightarrow K\)
is onto and has kernel \(\mathcal R\).
D.43.7 therefore gives
\[
 \begin{split}
 \pi(O^p(K_{\mathfrak F}))&=O^p(K),\\
 G&\simeq
 \mathfrak F/
       \bigl(O^p(K_{\mathfrak F})\mathcal R\bigr).
 \end{split}
 \tag{D.44.10}
\]
The product in the second line is closed:
it is the product of two compact normal
subgroups.
This proves the actual quotient comparison
for finite or infinite \(\Gamma\).

Choose the given conjugators of
\(\alpha_0(s)\) and \(\alpha_0(r)\)
in \(\Gamma\).
For the third pants generator in
\textup{(D.44.9)} choose the specific
conjugator
\[
                         h_3=\alpha_0(s)s .
 \tag{D.44.11}
\]
It works because \(\alpha_0(\tau)=\tau\):
conjugating \(s\tau s\) by
\(\alpha_0(s)s\) gives
\(\alpha_0(s)\tau\alpha_0(s)=\alpha_0(rs)\).
D.42, or D.43 for the infinite base,
provides a peripheral automorphism of
\(\mathfrak F/O^p(K_{\mathfrak F})\)
with these specified conjugator
projections.
It preserves the closed normal subgroup
generated by \(a^2,b^3\), since it sends
each \(a,b\) to a conjugate of itself.
The same argument for its inverse proves
equality, so it descends through
\textup{(D.44.10)} to an automorphism
\(\phi\) of \(G\).
It induces \(\alpha_0\).

Let \(A\in G\) be its third pants
conjugator, with projection
\(\alpha_0(s)s\).
Then
\[
 \begin{split}
 \phi(rs)&=A(rs)A^{-1},\\
 c&=\phi(s)A s\in\ker(G\to\Gamma),\\
 \phi(\tau)&=c\tau c^{-1}.
 \end{split}
 \tag{D.44.12}
\]
For the middle assertion the projection
is
\(\alpha_0(s)\alpha_0(s)ss=1\).
The last equation follows from
\(\tau=s(rs)s\).
Consequently
\(\phi_1=\operatorname{Inn}(c^{-1})\phi\)
fixes \(\tau\) and still induces exactly
the old \(\alpha_0\), not an unspecified
inner conjugate of it.

\paragraph{The exact abelianization and the kernel cusp gauge.}
Write \(R=\ker(G\to\Gamma)=K^{(p)}\)
and \(H=\ker(G\to\mathfrak S_3)\).
The source map \(F\to H\) is onto.
Let
\(\operatorname{ab}_p:F\to\mathbf Z_p^2\)
be the abelianization in the actual
basis \(x,y\).
Its restriction to \(K\) has pro-\(p\)
image, so it kills \(O^p(K)\).
It therefore defines a homomorphism
on \(H\).
Together with \(H\to\Gamma_F^{\mathrm{ab}}\)
it gives the exact isomorphism
\[
 H^{\mathrm{ab}}\ \xrightarrow{\ \simeq\ }\
       \mathbf Z_p^2\times\Gamma_F^{\mathrm{ab}}.
 \tag{D.44.13}
\]
Here \(\Gamma_F\), and thus its
abelianization, is pro-prime-to-\(p\).
To verify the assertion, every finite
abelian \(p\)-quotient of \(H\) is a
quotient of \(F^{\mathrm{ab},p}\);
the displayed map retains all of
\(\mathbf Z_p^2\), so it is the full
\(p\)-primary factor.
Every finite abelian quotient of
order prime to \(p\) kills \(R\)
and factors through \(\Gamma_F\).
These two assertions detect every
finite abelian quotient and prove
injectivity of \textup{(D.44.13)}.
The image is closed and surjects
onto each factor.
The two factors have no nontrivial
common finite quotient, since
their prime supports are disjoint.
The finite subdirect-product argument,
followed by compactness, proves
surjectivity.
This establishes the displayed
actual abelianization map.

Set \(S=\phi_1(s)\) and \(f_1=sS\).
Its projection is exactly \(f_0\).
Equation \textup{(D.44.2)} implies
that \(\alpha_0\) is the identity
on \(\mathfrak S_3\):
it fixes the image of \(s\), and
then that of \(r=s\tau\).
Thus \(\phi_1\) also induces the
identity there, and \(f_1\in H\).
Since \(S^2=1\),
\[
                       (s f_1s)f_1=1.
 \tag{D.44.14}
\]
The action of \(s\) on
\(\mathbf Z_p^2\) interchanges
the two coordinates by
\textup{(D.44.8)}.
Hence
\(\operatorname{ab}_p(f_1)=(a,-a)\)
for a definite \(a\in\mathbf Z_p\).
Its other coordinate in
\textup{(D.44.13)} is zero,
because \(f_0\) is in the image
of \(F'\).

Let \(\epsilon_p\in\widehat{\mathbf Z}\)
be the idempotent with \(p\)-component
one and every other primary
component zero.
The element
\[
                         g=x^{\epsilon_p a}
 \tag{D.44.15}
\]
centralizes \(\tau\) and belongs
to \(R\): every finite image of
\(\Gamma\) has order prime to \(p\),
so the image of this power there
is the identity.
Put \(\sigma=\operatorname{Inn}(g)\phi_1\).
Then \(\sigma\) still fixes \(\tau\)
and induces exactly \(\alpha_0\).
Its conjugator is
\[
 \begin{split}
 f=s\sigma(s)
   &=y^{\epsilon_p a}f_1x^{-\epsilon_p a},\\
 \operatorname{ab}_p(f)
   &=(a,-a)+(-a,a)=(0,0).
 \end{split}
 \tag{D.44.16}
\]
Its \(\Gamma_F^{\mathrm{ab}}\)
coordinate stays zero.
Thus \textup{(D.44.13)} proves
\(f\in H'\).
Surjectivity of \(F\to H\)
identifies \(H'\) with the
compact image of \(F'\):
the image contains all
commutators, and compactness
retains their closed subgroup.
This proves the first line
of \textup{(D.44.4)}, including
the unchanged literal \(f_0\).

\paragraph{The centralizer ambiguity is an exact kernel gauge.}
We record why the choice made
in \textup{(D.44.12)} introduces
no additional normalized
representative in the same
kernel-inner orbit.
First let \(\Gamma\) be finite
and put \(e=\operatorname{ord}_\Gamma(\tau)\).
The image of \(\tau\) in
\(\mathfrak S_3\) is a
transposition, so \(e\) is even.
The subgroup \(K\) corresponds
to the actual finite cover of
the three-punctured sphere
given by \(K\subset F\).
The word \(\tau^e=x^{-e/2}\)
is one of its based boundary
generators, with its orientation
possibly reversed.
There are \(|\Gamma|/e\geq3\)
boundary components: the coset
map \(\Gamma/\langle\tau\rangle
\to\mathfrak S_3/\langle\bar\tau\rangle\)
is onto.
Hence each boundary is a
primitive free generator after
an actual surface basis change,
as in D.41.
The Magnus centralizer comparison
there gives
\[
 C_R(\tau)\subset C_R(\tau^e)
   =\langle\tau^e\rangle
   =\langle\tau^{\epsilon_p}\rangle.
 \tag{D.44.17}
\]
The reverse inclusion is
immediate, and the last equality
uses that \(e\) is a unit in
\(\mathbf Z_p\).
Thus
\(C_R(\tau)=\langle x^{\epsilon_p}\rangle
\cong\mathbf Z_p\), since \(p\geq5\).

For an infinite \(\Gamma\),
take cofinal finite quotients
retaining \(\mathfrak S_3\).
The residual identity D.43.9,
whose proof applies equally
to the source \(\widehat B\),
identifies \(G\) with the
inverse limit of the corresponding
relative groups.
At every finite stage the
centralizer in its pro-\(p\)
kernel has the same distinguished
generator \(x^{\epsilon_p}\).
The transition maps carry
that generator to itself.
The parameters in
\(\mathbf Z_p\) are unique
by \textup{(D.44.17)}, so
any element of \(C_R(\tau)\)
has one common parameter.
This proves the same
centralizer equality in
the limit.

Two kernel conjugators
normalizing \(\phi(\tau)\)
differ by an element of
this centralizer.
Thus the resulting
\(\tau\)-fixed automorphisms
differ by
\(\operatorname{Inn}(x^{\epsilon_p v})\).
Formula \textup{(D.44.16)}
changes the first abelian
coordinate from \(a\) to \(a-v\).
The subsequent gauges
therefore cancel exactly:
\[
 \operatorname{Inn}(x^{\epsilon_p(a-v)})
 \operatorname{Inn}(x^{\epsilon_p v})
       =\operatorname{Inn}(x^{\epsilon_p a}).
 \tag{D.44.18}
\]
This proves uniqueness of
the representative fixing
\(\tau\) and having \(f\in H'\)
within the indicated kernel-inner
orbit.
Every normalization used here
is trivial on \(\Gamma\).

\paragraph{The two cycles, generation, and the actual symmetry.}
Let \(S=\sigma(s)\), \(R_0=\sigma(r)\).
Since \(S R_0=\tau\) and
\(f=sS\), we have the exact
formulas
\[
                  S=s f,\qquad R_0=f^{-1}r .
 \tag{D.44.19}
\]
The first follows from \(s^2=1\);
for the second use \(S^2=1\)
and \(sfs=f^{-1}\).
By \textup{(D.44.8)}, conjugation
by \(s\) evaluates a word
at \((y,x)\), and conjugation
by \(r\) evaluates it at
\((y,z)\).
Thus \(S^2=1\) gives the
first equation of
\textup{(D.44.5)}.
The inverse of \(R_0^3\) is
\[
       (r^2 f r^{-2})(r f r^{-1})f
              =f(z,x)f(y,z)f(x,y).
 \tag{D.44.20}
\]
This proves the second equation,
with its precise order.
Any preimage of \(f\) in \(F'\)
has these evaluation images,
because the evaluations in
\textup{(D.44.8)} are actual
conjugation maps on \(H\).
Also
\[
      \sigma(y)=SxS=f^{-1}yf,\qquad
      \sigma(z)=R_0^2 x R_0^{-2}.
 \tag{D.44.21}
\]
The second element is conjugate
to \(z=r^2xr^{-2}\) by
\(R_0^2r^{-2}\).
The restriction of \(\sigma\)
to \(H\) is an automorphism,
since it induces the identity
on \(G/H=\mathfrak S_3\).
This proves the required
generation in the actual
pure quotient.

For the symmetry comparison
let \(\iota_b\) denote
conjugation on \(H\) by
\(b\in G\).
Then
\[
 \sigma|_H\,\iota_b\,(\sigma|_H)^{-1}
       =\iota_{\sigma(b)}
       =\operatorname{Inn}(\sigma(b)b^{-1})\,\iota_b.
 \tag{D.44.22}
\]
The inner factor lies in \(H\),
because \(\sigma\) is the
identity on \(\mathfrak S_3\).
Passing to \(\operatorname{Out}(H)\)
proves the asserted centralizer
condition for the actual
geometric \(\mathfrak S_3\).

Equations \textup{(D.44.5)}
are precisely Harbater--Schneps
relations (I)--(II) at
\(\lambda=1\)
\cite[\S0.1, (I)--(II);
Proposition~4]{DE_HarbaterSchneps2000}.
Their three-strand comparison
also has the following literal
modular coordinates:
\[
                    u=r^{-1}s=\tau^{-1},
           \qquad v=sr^{-1}=s u s .
 \tag{D.44.23}
\]
Here \(uvu=vuv=s\) and \(uv=r\).
The automorphism just constructed
sends \(u\) to \(u\) and
\(v\) to \(f^{-1}v f\).
These equations follow from
\(\tau\)-fixing, \(sfs=f^{-1}\),
and \textup{(D.44.19)}.
Thus the actual source map
from the three-strand braid
presentation is retained.
Choose a cofinal sequence of characteristic open normal
subgroups of the finitely generated group \(G\), for example
the intersections of all open normal subgroups of index
at most \(n\). Intersect with \(H\), and, when \(\Gamma\)
is finite, also with \(R\); these intersections remain
\(\sigma\)-stable because \(\sigma\) preserves \(H,R\).
The resulting actual finite quotient maps commute with
\(\sigma\), retain the specified source generators, and
form a cofinal compatible system. In each of these
finite quotients the
torsion-to-hexagon calculation
of E.9, or the displayed
formulas themselves, gives
the complete two gentle
hexagons with multiplier one.
The map is onto because
\(\sigma\) is an automorphism
of \(G\).
The actual descent through the \(\sigma\)-stable kernels
supplies an automorphism at every stage.
One choice of a preimage
of \(f\) in \(F'\) works
simultaneously in all these
finite quotients.
This proves their compatible
gentle-shadow assertion.

\paragraph{Where the torsion gate is automatic and how the old braid shadow is retained.}
Suppose \(\Gamma_F\) is
pro-\(\ell\), \(\ell\geq5\),
and the old gentle
automorphism is already
normalized as in
\textup{(D.44.2)}.
It induces the identity on
\(\mathfrak S_3\).
In the inverse image of
\(\langle\bar s\rangle\),
the original \(\langle s\rangle\)
and the new
\(\langle\alpha_0(s)\rangle\)
are complements to the same
normal pro-\(\ell\) kernel.
The coprime complement proof
in D.41 conjugates them
by that kernel.
Their positive generators
have the same image
\(\bar s\), so the conjugacy
is on the generators
themselves.
Apply the same argument
to the inverse image of
\(\langle\bar r\rangle\).
Its quotient has order three,
and the images of the two
positive generators are
both \(\bar r\).
This proves exactly the
two conjugacies required
in the statement.

For the D.7--D.8 example,
quotient the actual three-strand
group by its central
generator.
The pure kernel over
\(\mathfrak S_3\) is
\[
                 F/(\gamma_{11}F\,F^{101}),
 \tag{D.44.24}
\]
a finite \(101\)-group.
The already proved gentle
maps are automorphisms,
and retain multiplier one
and their specified \(f_0\).
Hence the preceding complement
argument applies to every
one of them.

To retain also the original
finite braid target
\(\widetilde\Gamma=B_3/K_3\),
form the actual common
refinement
\[
             G\mathbin{\times_\Gamma}
                         \widetilde\Gamma .
 \tag{D.44.25}
\]
The kernel of \(G\to\Gamma\)
is pro-\(p\), and the
kernel of
\(\widetilde\Gamma\to\Gamma\)
is the central group of
order \(101\).
The natural source map
\(\widehat B_3\) to
\textup{(D.44.25)} is onto:
on finite kernel quotients
its two Goursat quotients
have orders dividing
both a power of \(p\)
and \(101\), so they
are trivial; compactness
then gives surjectivity.
The two automorphisms
\(\sigma\) and the old
braid automorphism induce
the same \(\alpha_0\)
on \(\Gamma\).
Their product is therefore
an actual automorphism
of the displayed fibre
product.

Finally the direct product
\(PB_3=F\times\langle C_3\rangle\)
used in D.7 shows that
the image of \(F\) in
\(\widetilde\Gamma\) maps
isomorphically to the
pure group
\textup{(D.44.24)}.
Thus any preimage in \(F'\)
of the new \(f\), already
projecting to \(f_0\) in
\(\Gamma\), has exactly
the old \(f_0\) in
\(\widetilde\Gamma\);
there is no discarded
central parameter in
this comparison.
On \textup{(D.44.25)}
the braid action
\(u\mapsto u,\ v\mapsto f^{-1}vf\)
agrees with the two
component automorphisms.
It fixes the old central
generator because the
multiplier is one.
This proves the full
old-shadow assertion,
including its central
coordinate, generation,
and all projections.
\end{proof}

The program
\texttt{a\_modular\_coprime.py}
retains all twelve literal
Schreier expressions in
\textup{(D.44.7)}, the
complete free-product
identities
\textup{(D.44.8)--(D.44.9)},
and both cycle comparisons.
It independently checks
the signed integer cusp
gauges.
Its finite model is the
actual quotient of
\(C_2*C_3\) with pure
kernel \(\mathbf F_5^2\)
and the displayed
\(\mathfrak S_3\) action,
of order \(150\).
All \(25\) kernel-inner
maps and all five choices
of cusp centralizer
give the same normalized
representative; all
\(125\) cases, including
their unique scalar
gauges, are saved.

The conclusion concerns
the complete coprime
relative covers of the
specified old gentle
automorphisms.
Finite order-two and
order-three equations
alone do not supply the
input conjugacies in a
general mixed-primary
quotient; the proof
establishes them for
the pro-\(\ell\) inputs
just specified.
The remaining D.2
comparison still includes
relative primes dividing
the old base and the
other finite source
extensions.
